\begin{lemma}
\label{editorial-source-lemma-geometrically-connected-any-base-change}
Let $k$ be a field.
Let $S$ be a geometrically connected $k$-algebra.
Let $R$ be any $k$-algebra.
The map
$$
R \longrightarrow R \otimes_k S
$$
induces a bijection on idempotents, and the map
$$
\Spec(R \otimes_k S) \longrightarrow \Spec(R)
$$
induces a bijection on connected components.
\end{lemma}

\begin{proof}
Put $A=R\otimes_k S$ and retain
$f:X=\Spec(A)\to Y=\Spec(R)$.
Geometric connectedness of the
original $S$ includes nonemptiness,
so $S\neq0$. If $R=0$, both
spectra are empty and both
rings have one idempotent,
so the two asserted bijections
hold. Assume $R\neq0$.
For every $\mathfrak p\in Y$,
the original fibre comparison
$$
A\otimes_R\kappa(\mathfrak p)\longrightarrow
\kappa(\mathfrak p)\otimes_k S,
\qquad(r\otimes s)\otimes a\longmapsto ra\otimes s
$$
has inverse $a\otimes s\mapsto(1\otimes s)\otimes a$.
The receiving spectrum is
connected by the hypothesis
on $S$ and is nonempty.
Hence $f$ is surjective with
connected fibres. It is open
without finite type assumptions,
by Lemma
\ref{editorial-source-lemma-map-into-tensor-algebra-open}.

Let $C\subseteq X$ be open
and closed. Connectedness
of each fibre implies that
the fibre is contained either
in $C$ or in its complement:
otherwise their intersections
would separate the fibre.
The images $f(C)$ and $f(X\setminus C)$
are therefore disjoint open
sets covering $Y$. In
particular $f(C)$ is open
and closed, and
$$
C=f^{-1}(f(C)).
$$
Conversely, inverse images of
open and closed subsets are
open and closed, and surjectivity
gives $f(f^{-1}(D))=D$ for
every such $D\subseteq Y$.
This proves a bijection of
the original open-and-closed
subsets in both directions.

By Lemma \ref{algebra-lemma-disjoint-decomposition},
the map $e\mapsto D(e)$ is
a bijection between idempotents
of each ring and its open-and-closed
subsets. Under the original
ring map $R\to A$, the
identity $f^{-1}(D(e))=D(e\otimes1)$
therefore proves exactly the
claimed idempotent bijection.
The inverse takes an idempotent
$a\in A$ to the unique
idempotent $e\in R$ such
that $D(e)=f(D(a))$.
These maps also preserve
complements and intersections,
represented by $1-e$ and
$ee'$, so give an isomorphism
of the corresponding Boolean
algebras. They commute with
arbitrary original $k$-algebra
base change: the image of
$e\otimes1$ is $\theta(e)\otimes1$,
and uniqueness of descent
gives the inverse compatibility.

We prove the component assertion
directly as well. For any
connected component $D$ of
$Y$, the restriction
$f^{-1}(D)\to D$ is open,
surjective and has the same
connected fibres. Openness
holds because for every open
$U\subseteq X$,
$$
f(U\cap f^{-1}(D))=f(U)\cap D.
$$
If $f^{-1}(D)$ were separated
into disjoint nonempty open
and closed parts, the same
fibre argument would give
a separation of $D$ by
their open images. Thus
$f^{-1}(D)$ is connected
and nonempty. It is maximal
connected: any larger connected
subset would have connected
image contained in a component
of $Y$, which must be $D$
because it meets $D$.
Conversely each component
of $X$ has connected image
in some component $D$ of
$Y$ and hence equals this
connected inverse image.
This proves the claimed
bijection and the exact
inverse-image description,
as in Topology, Lemma
\ref{topology-lemma-connected-fibres-connected-components}.

With the quotient topologies
on the sets of components,
the bijection is a homeomorphism.
Write $q_X,q_Y$ for the two
quotient maps and $h$ for
the induced bijection.
The relation $h q_X=q_Y f$
proves continuity by the
quotient property. If $V$
is open in the component
space of $X$, then
$$
q_Y^{-1}(h(V))=f(q_X^{-1}(V))
$$
by the exact component
correspondence. The right
side is open because $f$
is open, proving that $h$
is open and hence a
homeomorphism. None of this
requires individual components
to be open in their spectra.
\end{proof}